\section{Introduction}
\label{sec:introduction}

\begin{table*}[t]
\centering
\footnotesize
\caption{Comparison of \hsfp{}~\cite{tran2026h} and \ehsfp{}. Both use class-wise statistics and local training within a client--edge--cloud hierarchy.}
\label{tab:extension_comparison}
\setlength{\tabcolsep}{3pt}
\renewcommand{\arraystretch}{1.08}
\begin{tabular}{@{}p{0.16\textwidth}p{0.31\textwidth}p{0.48\textwidth}@{}}
\toprule
Aspect & \hsfp{} & \ehsfp{} \\
\midrule
Statistical inputs
& Received class-wise means and coordinate-wise standard deviations.
& Timestamped moment replay, source masking, optional accumulation, and learned source weighting. \\
\addlinespace
Synthesis and training
& Diagonal Gaussian synthesis with local contrastive and task objectives.
& Prototype replay consistency and optional full covariance shared across classes for classification. \\
\addlinespace
Analysis
& Global-classifier gradient norms averaged over communication rounds.
& Operator-level characterization of replay, masking, and aggregation; conditional tier-wise stationarity relative to fresh-prototype base objectives, with reference variation, drift, additional update errors, and regularization. \\
\addlinespace
Evaluation
& Classification, segmentation, and sensitivity studies.
& Source-mask and delay sweeps, simulated execution disruptions, and replay, aggregation, and synthesis configuration contrasts. \\
\bottomrule
\end{tabular}
\end{table*}

Federated learning (FL) trains a shared model while raw examples remain on participating clients. FedAvg~\cite{mcmahan2017communication} and FedProx~\cite{li2020federated} aggregate locally trained models, so each client must train a full copy of the model. Split learning (SL) instead distributes layers between clients and servers to reduce client-side computation~\cite{vepakomma2018split}, and hierarchical FL inserts edge servers between clients and the cloud~\cite{liu2020client}, providing a further tier on which model segments can run. Conventional split training, however, exchanges per-sample activations and gradients at every cut, so its communication grows with the amount of data processed.

Prototype-based FL and server-side distillation coordinate learning through class representations or transferred knowledge~\cite{tan2022fedproto,lin2020ensemble,zhu2021data}. Hierarchical split-federated prototyping (\hsfp{})~\cite{tran2026h} brings this idea to a partitioned client--edge--cloud model. Clients report class-wise feature means, coordinate-wise standard deviations, and support counts. Edge and cloud learners then train their segments on synthetic features drawn from diagonal Gaussians, without cross-tier backpropagation. Communication thus scales with the number of classes rather than the number of examples. The price of this decoupling is that, apart from the separate ISIC-2018 decoder stage, the upper tiers never see real features: their training distribution is whatever the incoming summaries define.

\noindent\textbf{Research gap.}~
Partial participation and packet loss alter the statistics of classes that remain represented and remove other classes altogether. Delayed packets describe features produced by earlier encoder states. Asynchronous FL analyses model delay as staleness of parameters or gradients~\cite{xie2019asynchronous,nguyen2022federated,fraboni2023general}. At a statistical interface, by contrast, availability and age change the distribution from which the learner samples its inputs, and the learner has no access to the distribution it should have trained on. \hsfp{} discusses decoupled asynchronous training, but its convergence bound leaves the effects of incomplete and stale summaries implicit. Two questions therefore remain open: how each mechanism that handles intermittent summaries changes the synthetic training distribution, and how that change enters the optimization guarantee of an upper-tier learner.

\noindent\textbf{Approach.}~
We propose Episodic Hierarchical Split-Federated Prototyping (\ehsfp{}). Event-driven FL systems motivate separating individual processing phases from persistent learning state~\cite{grafberger2021fedless,qi2024lifl}. In our simulator, an episode is one processing phase at an edge or cloud learner. Each episode constructs synthetic training inputs from the available statistics, while the model, optimizer, and replay state persist across phases. A timestamped memory stores class moments, and moment replay mixes stored and received moments before synthesis. Training-time source masking varies which source outputs contribute after replay. A learned scorer weights the remaining outputs using class support, feature spread, packet age, and the distance between a source's mean and the average source mean for that class. Prototype replay consistency (PRC) penalizes differences between the learner's outputs on independent draws from replayed and stored summaries. For classification, the generator can also use a full covariance shared across classes. \ehsfp{} retains the hierarchy, statistical interface, local objectives, and separate edge and cloud synchronization schedules of \hsfp{}~\cite{tran2026h}, the closest prior framework. Table~\ref{tab:extension_comparison} summarizes the differences.

\noindent\textbf{Analysis.}~
We write the synthesis pipeline as a composition of replay, masking, accumulation, aggregation, and sampling operators. We measure its effect by the 2-Wasserstein drift between the actual synthetic law and a reference built from complete, fresh summaries. Replay adds drift at most linearly in the weight placed on memory when current summaries are fresh, and it restores coverage at a cost that grows with record age when they are not. Masking removes a class with a probability determined by how many sources hold it. Learned weighting can move the synthetic mean away from the pooled mean only when sources disagree. For $T$ valid SGD-type updates of one upper-tier learner, the optimization and stochastic-noise terms of our stationarity bound scale as $O((1+V_T)/\sqrt{T})$, where $V_T$ bounds the total reference variation. The bound retains residual terms for drift, additional mean update errors, and regularization.

\noindent\textbf{Evaluation.}~
We evaluate \ehsfp{} on four classification datasets and on ISIC-2018 segmentation, with five runs of 200 rounds per setting. Mean CIFAR-100 accuracy is 56.07\%, and the 50-client ISIC-2018 pipeline reaches 69.12\% IoU and 80.23\% Dice. Baseline values are taken from~\cite{tran2026h} (Section~\ref{sec:setup}). Under $\mathrm{Dir}(0.3)$, \ehsfp{} attains 9.13\% accuracy on CIFAR-100, so severe statistical heterogeneity remains difficult.

\noindent\textbf{Contributions.}~
\begin{enumerate}
    \item \emph{Formulation.} We formulate upper-tier learning as optimization over synthetic distributions produced by an explicit operator pipeline, which separates prototype drift, class coverage, the weighting-induced reference gap, and the Gaussian representation error (Sections~\ref{sec:preliminary} and~\ref{sec:formulation}).
    \item \emph{Mechanism analysis.} We prove that the statistical interface pools exactly under count weights and derive closed-form expressions for the drift added by replay, the class loss caused by masking, and the mean shift introduced by learned weighting (Lemma~\ref{lem:exact_pooling} and Propositions~\ref{prop:replay_drift}--\ref{prop:weighting_shift}).
    \item \emph{Method.} We develop \ehsfp{}, which adds timestamped moment replay, learned source weighting, source masking, PRC, and optional shared-covariance synthesis to the \hsfp{} interface (Section~\ref{sec:method}).
    \item \emph{Optimization guarantee.} We bound the average squared gradient norm of the fresh-prototype base objective of one upper-tier learner, with drift entering as an explicit residual, and bound that drift by contribution age under additional assumptions (Theorem~\ref{thm:convergence} and Lemma~\ref{prop:drift_decomp}, Section~\ref{sec:convergence}).
    \item \emph{Evaluation.} We report classification and segmentation results, source-mask, delay, and execution-event sweeps, and replay, aggregation, and synthesis configuration contrasts (Section~\ref{sec:experiments}).
\end{enumerate}
